% Replace the Resource-quality audit paragraph in Section 8 (Evaluation Protocol).
% Planned protocol: keep future tense until the audit has actually been completed.
\subsection{Human Annotation Audit}
\label{sec:human-annotation-audit}

We will evaluate annotation quality through a sampled human audit. Three
reviewers will independently assess the same 140 question--answer records,
together with their supporting structured and language annotations. A separate
set of 20 shared calibration examples will be used to refine the rubric and
excluded from formal analysis. Each reviewer will therefore assess 160
examples, yielding 420 formal judgments and 60 calibration judgments.

\paragraph{Audit endpoint.}
The primary judgment will be binary: acceptable or containing a material
error. Reviewers will check question validity, reference-answer consistency
with the verification record, language fidelity, evidence support, and
compliance with the task's public/private information boundary. Criteria will
be applied according to question type; a clarification item will be assessed
for whether the missing information and requested clarification are justified.
Reviewers may inspect private verification evidence for auditing, although
these fields remain unavailable to evaluated answering models.

\paragraph{Sample-size justification.}
We adopt an 80\% nominal confidence level, corresponding to a two-sided
significance level of $\alpha=0.20$, and target 80\% statistical power.
For three raters, the binary multi-rater agreement calculation implemented
in \texttt{kappaSize::PowerBinary} \citep{rotondiKappaSize} compares
$H_0:\kappa=0.60$ with an anticipated alternative of $\kappa=0.80$,
assuming an acceptable-rating prevalence of $p=0.90$.
These inputs yield an approximate requirement of 129.03 records, rounded up
to 130. We select 140 formal records, providing approximately 82.4\% modeled
power under the same assumptions. The 90\% prevalence is a planning
assumption, rather than an observed acceptance rate or a dataset acceptance
threshold. This calculation concerns pooled reviewer agreement; it does not
power separate domain--question-type comparisons or establish a lower bound
of 90\% on dataset correctness.

\paragraph{Sampling and reporting.}
The sample will cover the released domains and question types. Sampling will
avoid repeated records from the same trajectory to support the independence
assumption. The sampling seed, allocation, rubric, and reviewer qualifications
will be documented with the results. We will report multi-rater agreement
before adjudication, raw agreement, and per-reviewer acceptance rates.
Disagreements will subsequently be adjudicated against verification evidence,
with error categories, resulting corrections, and the adjudicated acceptance
rate reported separately. Actual uncertainty estimates will accompany the
observed results; modeled power remains conditional on the assumed rating
distribution and agreement model.
